% Path Complexity Metrics - LaTeX Table for Publication
% Auto-generated by Step 6

\begin{table}[htbp]
\centering
\caption{\textbf{Path complexity metrics for age-group trajectories (2007--2024).} 
Tortuosity measures the ratio of actual path length to straight-line displacement; 
values near 1.0 indicate direct movement while higher values reflect meandering or 
period-specific fluctuations. The mean resultant length (R) quantifies directional 
consistency, with values near 1.0 indicating all period-to-period transitions point 
in similar directions. Circular standard deviation provides an intuitive measure of 
angular spread.}
\label{tab:path_complexity}

\begin{tabular}{lccccc}
\toprule
\textbf{Age Group} & 
\textbf{Displacement} & 
\textbf{Path Length} & 
\textbf{Tortuosity} & 
\textbf{R} & 
\textbf{Circular SD} \\
 & (km) & (km) & (ratio) & & (degrees) \\
\midrule
10-17 & 0.810 & 1.832 & 2.26 & 0.169 & 108.0 \\
18-30 & 0.633 & 1.004 & 1.58 & 0.204 & 102.1 \\
31-55 & 0.502 & 0.510 & 1.02 & 0.983 & 10.7 \\
56-65 & 1.766 & 2.146 & 1.22 & 0.599 & 58.0 \\
66+ & 0.979 & 8.661 & 8.85 & 0.200 & 102.8 \\
\bottomrule
\end{tabular}

\vspace{0.3cm}
\begin{tablenotes}
\small
\item \textit{Displacement}: Straight-line distance from the first period (2007--09) to the final period (2022--24), measured on the same chain as the path length. Both coordinates enter unscaled, the horizontal as a fraction of the day and the vertical in km, so the ratio is not invariant to that choice of units.
\item \textit{Path Length}: Total distance traveled through action-space across all period-to-period transitions.
\item \textit{Tortuosity}: Path length / displacement. Values of 1.0 indicate perfectly straight trajectories; 
higher values indicate deviation from direct movement.
\item \textit{R (Mean Resultant Length)}: Circular statistic measuring directional consistency (range: 0--1). 
Higher values indicate period-to-period vectors point in more similar directions.
\item \textit{Circular SD}: Angular dispersion of trajectory directions in degrees. 
Lower values indicate tighter angular clustering.
\end{tablenotes}
\end{table}

